\section{Recovered grading, homogeneous algebra, and sharper type bounds}
\label{sec:homogeneous-integration}

The incoming package \path{unknot_recognition_dense_algebra_20261008.zip}
contains another dense coefficient engine, a recovered-grading argument,
and a bounded finite-quotient certificate search. It is now preserved as
report 16. Its component contraction
overlaps the implementation already integrated from reports 10 and 12.
We therefore retain the current adapter and incorporate the additional
homogeneous multiplication case, rather than replacing the production
tree with the archive's older baseline. We also derive a packed shortcut
for products in the top degree and audit the recovered grading against
the current scanner.

\subsection{The grading that survives after its labels are erased}
Start the empty object with an auxiliary quantum shift $q=0$.
When smoothing $i\in\{0,1\}$ of one crossing creates $c$ closed
circles, its delooping label $\ell$ gives the child shift
\[
 q'=q+i+c-2\operatorname{popcount}(\ell).
\]
For matching objects $a_u,a_v$ at a frontier of $2m$ endpoints, every
monomial with $d$ dots in a nonzero differential entry $u\to v$ obeys
\begin{equation}\label{eq:recovered-grading}
 c(a_u,a_v)-m-2d+q_v-q_u=0.
\end{equation}
Here $c(a_u,a_v)$ counts the circles of their overlay, and the
normalization is that of the raw scanner, before oriented link shifts.
This is the dotted grading of
\href{https://arxiv.org/abs/math/0410495v2}{Bar--Natan's tangle formalism},
specialized to the implementation's smoothing and delooping order.

\begin{proposition}
Every actual differential entry in the ordinary scanner is homogeneous
in dot degree. Every cancellable entry between equal matchings is the
scalar identity $1$.
\end{proposition}
\begin{proof}
A $k$-disc cobordism carrying $d$ dots has normalized degree $k-m-2d$.
A crossing saddle contributes $-1$, compensated by the smoothing
shift. Delooping a circle contributes shifts $+1$ and $-1$, giving
the displayed recurrence and establishing~\eqref{eq:recovered-grading}
after crossing extension.

The intrinsic weights from Section~\ref{sec:radical-adaptive} add under
nonzero composition; their negatives are these normalized degrees.
If a map between equal matchings has scalar coefficient one,
homogeneity forces its dot degree to be zero. There is only one
zero-dot monomial, so the map is exactly $1$ and the two shifts agree.
The source and target shifts telescope in a Schur correction through
this identity. That correction has the same degree as the existing
entry, and their sum remains homogeneous. Induction over the
cancellation sequence proves the claim.
\end{proof}

The auxiliary labels are a proof and diagnostic device; the normal
scanner still does not store quantum grading. Component sharing
preserves the conclusion: identical connected keys have identical
relative shift differences along their nonzero entries, and an
overall shift of a disconnected summand is irrelevant to homogeneity.
The adaptive degree-gap shortcut produces a zero differential, which
also admits such a grading before subsequent crossing extension.
This does not turn its output into a computation of quantum-graded
homology; those discarded absolute labels are not reconstructed.

The production audit explicitly propagates $q$ and checks every term
before and after elimination, for ordinary and forced component
composition. It completed 520 scans of 87 diagrams, using input,
greedy, and seeded shuffled orders where specified: 196,928 entries
and 220,736 supported terms satisfy the equation. The largest dot
degree observed was two. This finite observation is not a universal
bound on dot degree. Inputs, orders, exact ranks, and counts are in
\path{fast/results/grading_audit_20261008.json}; reproduce with
\path{fast/audit_grading.py}.

\subsection{Cheaper multiplication with an explicit applicability check}
In the component ring
$R_r=\F[x_1,\ldots,x_r]/(x_1^2,\ldots,x_r^2)$, suppose $F,G$ have
homogeneous dot degrees $a,b$. Apply the ordinary subset zeta
transform to their binary coefficient arrays, multiply pointwise,
and apply the inverse transform, also XOR over $\F$. The resulting
coefficient at $S$ is
\[
 U_S=\sum_{A\cup B=S}F_A G_B.
\]
Retain only subsets of size $a+b$. Each supported pair has
$|A|+|B|=a+b$, so this restriction is equivalent to $A\cap B=\varnothing$.
The surviving array is exactly the square-free product $FG$.
If $a+b>r$, the product is zero.

This branch uses $O((r+1)2^r)$ binary field operations. The generic
ranked transform, retained for mixed-degree inputs, has the usual
$O((r+1)^2 2^r)$ ring-operation bound of
\href{https://arxiv.org/abs/cs/0611101v1}{Bjoerklund--Husfeldt--Kaski--Koivisto}.
These are arithmetic bounds; array indices, monomial masks, and
coefficient conversions still carry their bit costs. Both remain
within the earlier $\operatorname{poly}(r)2^r$ coefficient envelope.
The ranked implementation itself is cheaper on many special inputs,
so the two displayed bounds are not a prediction of a factor-$r$
measured speedup on every homogeneous product.

The final size restriction is essential. For example,
$x_1(x_1+x_2)=x_1x_2$, whereas unfiltered union convolution would
incorrectly keep $x_1$. The software verifies the dot degree of each
actual operand and falls back to ranked convolution for mixed degrees.
It never infers this property merely from a caller selecting the
scanner API. Projection to component variables preserves degree for
every surviving monomial; a repeated component annihilates a term
and collisions may cancel terms, but cannot create another degree.
Thus the optimization applies after the existing component projection.

\paragraph{Complementary subsets in the top degree.}
We add a further specialization when $a+b=r$. Only the full monomial
$x_{[r]}$ can survive, and its coefficient is
\[
 [x_{[r]}]FG=\sum_{A\subseteq[r]}F_A G_{[r]\setminus A}.
\]
In the packed coefficient convention, complementing the subset index
is exactly reversing its $2^r$ coefficient bits. Reverse the bytes
and the bits within each byte, remove padding when $2^r<8$, intersect
with the other polynomial, and take the parity of the population
count. This computes the coefficient using linear work in the packed
bit length after the degrees have been checked, with no zeta
transform. The homogeneity check still incurs its own support and
index-processing cost. Generic mixed-degree inputs cannot use this
shortcut to compute their other output degrees.

The new cases are inside the existing opt-in component-contraction
backend. No additional recognition flag is required; the ordinary
sparse engine remains the default. The old ranked function remains
available for independent comparisons. Existing allocation limits,
deadline callbacks, scalar identities, and polarization are retained.

\subsection{A sharper sufficient regime for component sharing}
Let $B_m=\binom{m}{\lfloor m/2\rfloor}$ and let
$M(2m)=(2m-1)!!$, with $M(0)=1$. The latter is the safe unrestricted
pairing count; this argument does not restore a Catalan bound without
a common-disk certificate.

A homogeneous coefficient with at most $m$ circle variables lies
in one of at most $m+1$ degree layers, each of dimension at most
$B_m$. There are therefore at most
\[
 H_m=(m+1)2^{B_m}
\]
possibilities, intentionally counting zero more than once. A connected
component on $s$ objects has normalized homological degrees between
zero and $s-1$: each edge changes degree by one, and a simple path
from a minimum-degree object has at most $s-1$ edges. Consequently
the number of ordered serialized keys with at most $b$ objects is
bounded by
\[
 K(m,b)=\sum_{s=1}^b s^s M(2m)^s H_m^{s^2}.
\]
We have overcounted zero entries, invalid degree differences,
disconnected graphs, and failures of $d^2=0$, so no unproved
realizability assertion enters this upper bound. Taking logarithms,
\[
 \log K=O\bigl(b^2B_m+bm\log(m+1)+b\log(b+1)\bigr)
       =O(b^2B_m).
\]

The existing exact-sharing argument now gives the stronger sufficient
condition
\[
 b^2\binom{m}{\lfloor m/2\rfloor}=O(\log^2(n+2))
\]
for quasi-polynomial cost, provided it bounds every reduced component
and frontier actually encountered. Extending a representative through
one crossing only multiplies its object count by a constant; different
parent components stay independent. Shift weights have polynomial
degree and polynomial bit length, as in Section~\ref{sec:structural}.
For bounded $b$, writing $L=\log_2(n+2)$ and $\ell=\log_2 L$, the
sufficient frontier scale improves from $2m\le4\ell+O(1)$ to
\[
 2m\le4\ell+\log_2\ell+O(1).
\]
This is a refinement of a conditional parameter regime, not a proof
that arbitrary diagrams admit such scans. Actual component size is
still uncontrolled, and large multiplicities still defeat an
explicit-object implementation.

\subsection{Adaptive controls that satisfy the grading constraint}
The mixed-dot two-term matrices used in the previous adaptive
benchmark are valid ungraded complexes, but generally cannot be
genuine scan differentials: their entries mix dot degrees. Their
large speedups therefore do not establish the behavior of a graded
scan. We retain those samples and add a separate graded family.

For size $s$, begin with an invertible binary $s$-by-$s$ matrix,
constructed by reversible row additions, and add one source column.
Interpret each nonzero entry as the identity of a fixed matching.
The resulting two-term complex has $s+1$ sources, $s$ targets,
quantum shifts all zero, and exactly one surviving source object.
It satisfies both $d^2=0$ and the recovered grading equation.
Dense scalar Schur updates can still be expensive even though they
require no cobordism products. The existing adaptive allowance detects
that update work and switches to binary residue homology.

The new graded controls are still synthetic: we do not claim that
our crossing order naturally generates their density. This makes
them a stronger applicability check for the switch, not evidence of
a broad recognition improvement.

\begin{center}
\begin{tabular}{@{}lrr@{}}
\toprule
Kernel & Product ratio & Full component ratio \\
\midrule
homogeneous\_8\_2\_3 & 1.14 & 1.08 \\
homogeneous\_10\_3\_4 & 1.08 & 1.07 \\
homogeneous\_12\_4\_5 & 1.19 & 1.13 \\
homogeneous\_14\_5\_6 & 1.20 & 1.13 \\
homogeneous\_8\_4\_4 & 16.25 & 4.69 \\
homogeneous\_10\_5\_5 & 22.62 & 5.21 \\
homogeneous\_12\_6\_6 & 28.86 & 5.97 \\
homogeneous\_14\_7\_7 & 34.21 & 6.54 \\
mixed\_12 & 1.01 & 1.00 \\
sparse\_12 & 1.10 & 1.11 \\
\bottomrule
\end{tabular}
\end{center}
\begin{center}
\begin{tabular}{@{}lrr@{}}
\toprule
Graded two-term objects & Adaptive ratio & A/A \\
\midrule
65 & 2.72 & 1.026 \\
129 & 8.91 & 1.005 \\
257 & 20.74 & 1.001 \\
513 & 58.52 & 0.987 \\
\bottomrule
\end{tabular}
\end{center}
\begin{center}
\begin{tabular}{@{}lrr@{}}
\toprule
Forced component scan & Ranked/new ratio & A/A \\
\midrule
trefoil & 0.803 & 1.058 \\
conway & 0.991 & 1.015 \\
kinoshita\_terasaka & 1.092 & 1.060 \\
torus\_3\_5 & 0.987 & 0.990 \\
hard\_unknot\_8 & 1.006 & 1.002 \\
\bottomrule
\end{tabular}
\end{center}
All speedups are median paired old/new times; values above one are faster.


All measurements use seven shuffled paired trials and an identical
old/old control. Product timings compare the prior ranked transform
to the checked dispatch; full-component timings include projection
and output expansion. Ordinary scan measurements force component
contraction to exercise the changed path and include order selection.
The small trefoil scan regressed in this sample; other scan differences
are modest and include control variation. The default sparse engine
does not use this new dispatch, and no general recognition speedup is
claimed. Reproduction scripts are \path{fast/benchmark_homogeneous.py}
and \path{fast/benchmark_adaptive_graded.py}; their complete samples
are in \path{fast/results/homogeneous_integration_20261008.json} and
\path{fast/results/adaptive_graded_20261008.json}.

The full suite passes 229 tests. New tests exhaust every pair of
four-variable homogeneous polynomials, test overlap rejection and
byte padding, compare mixed inputs with the ranked fallback, verify
projection behavior, and combine forced component algebra with adaptive
cancellation. The isolated incoming archive passed its 74 tests,
the independent eleven-test algebra suite, 2,100 associativity
checks, 8,768 transfer comparisons, and replay of ten finite-quotient
certificates. Local logs are retained under \path{synthesis/data/}.

\subsection{A limit of fixed finite-quotient portfolios}
The archive's optional $A_5$ search has not been merged into the
current pipeline. Its witnesses are useful nontriviality certificates;
exhaustion must remain inconclusive. The report also supplies a
structural limitation worth retaining in the research plan.

For prime $p>3$, the torus-knot group of $T(3,p)$ has presentation
$\langle x,y\mid x^3=y^p\rangle$, with $x^3=y^p$ central.
In any finite image in which the order of the image of $y$ is
coprime to $p$, that image is a power of its central $p$th power.
It is therefore central, so the image generated by $x,y$ is abelian.
In particular this holds for $A_5$ when $p\ge7$.

For a coloring by a finite quandle $Q$ with $|Q|<p$, its inner action
lies in the symmetric group on fewer than $p$ points, so the same
argument makes the induced group image abelian. All meridian images
then agree. The common right translation fixes each color by
idempotence, making every crossing preserve its under-arc color;
following the one-component knot forces the coloring to be constant.
Thus this family needs palettes of at least $p$ for nonconstant
colorings, although its usual braid diagram has only $2p$ crossings.
These knots are already decided by our three-braid or homogeneous
diagram stages. The obstruction rules out a universal fixed-palette
strategy; it does not rule out larger adaptive targets, bounded-seed
searches, or a quasi-polynomial algorithm by another method.
